% TOP_PROBLEM: {"id": 3700014, "title": "Composite periodic entire functions", "category_id": 8, "category": {"id": 8, "name": "analysis", "display_name": "Analysis", "description": "Limits, continuity, calculus, and function theory.", "slug": "analysis", "order_index": 8, "created_at": "2026-07-31T15:26:25.671Z"}, "status": "open", "problem_number": "AMR-036-0014", "set_id": 13, "statusQualification": "Excludes trivial constant compositions and clarifies that the quadratic alternative is a classification condition on an already periodic composition."}
% FIRST_POSTED: 2026-10-01
% ENTRY_KIND: statement_only
% UnsolvedMath ID 3700014; source catalogue code AMR-036-0014
% !TeX program = lualatex
% Definition and sources only; this file is not an LLM solution attempt.
\documentclass[11pt,a4paper]{article}
\usepackage[margin=25mm]{geometry}
\usepackage{fontspec}
\setmainfont{lmroman10-regular.otf}[BoldFont=lmroman10-bold.otf,
 ItalicFont=lmroman10-italic.otf,BoldItalicFont=lmroman10-bolditalic.otf]
\usepackage{amsmath,amssymb,mathtools}
\usepackage{unicode-math}
\setmathfont{latinmodern-math.otf}
\AtBeginDocument{\let\setminus\smallsetminus}
\usepackage{xurl,hyperref}
\hypersetup{colorlinks=true,urlcolor=blue}
\setcounter{secnumdepth}{0}
\setlength{\parindent}{0pt}
\setlength{\parskip}{5pt}
\setlength{\emergencystretch}{3em}
\begin{document}
\section*{Composite periodic entire functions}\label{3700014}
{\small UnsolvedMath ID 3700014 · AMR-036-0014 · Analysis · AMR Open Problem Lists}
\subsection{Definitions and mathematical statement}
Let $f,g:\mathbb C\to\mathbb C$ be nonconstant entire
functions. A nonzero number $T$ is a period of $f\circ g$
if $f(g(z+T))=f(g(z))$ for every $z\in\mathbb C$.
Classify the pairs for which such a period exists.

The proposed classification says that every pair belongs,
after affine changes of the intermediate variable, to at
least one of the following types. The function $g$ is
periodic; or there are $T,K\ne0$ with
$g(z+T)=g(z)+K$ and $f(w+K)=f(w)$; or $g$ is a
polynomial of degree two; or
\[
 g=P\circ h,\qquad \deg P=2,\qquad
 h(z+T)=h(z)+K,\qquad
 (f\circ P)(w+K)=(f\circ P)(w),
\]
where $h$ is entire and $T,K\ne0$.
In considering a type, one may choose a suitable nonzero
period, rather than fixing one particular period in advance.

An affine change of the intermediate variable means replacing
$(f,g)$ by $(f\circ L^{-1},L\circ g)$ with
$L(w)=aw+b$, $a\ne0$, which leaves the composition unchanged.
The quadratic alternative is a necessary structural
possibility for an already periodic composition, not the
assertion that every outer function produces a periodic
composition with every quadratic polynomial.
Prove that these mechanisms exhaust all possibilities,
or construct a pair outside them. Excluding constant
functions removes compositions that are trivially periodic
regardless of the inner function.
\subsection{Short English statement}
Must every periodic composition of nonconstant entire functions arise from the listed periodic, translation, or quadratic mechanisms?
\subsection{Sources}
\begin{itemize}
\item[S1] ULAM.AI and the UnsolvedMath Contributors, supplied problems.json, record 3700014 (AMR-036-0014), AMR Open Problem Lists. Catalogue identity and source transcription; CC BY 4.0.
\url{https://huggingface.co/datasets/ulamai/UnsolvedMath}.
\item[S2] Eremenko - My favorite unsolved problems, Composite periodic entire functions conjecture. Original source indexed by AMR-036-0014.
\url{https://www.math.purdue.edu/~eremenko/uns1.html}.
\item[S3] A. Eremenko, Composite periodic entire functions. Author-maintained primary problem note.
\url{https://www.math.purdue.edu/~eremenko/dvi/periodic.pdf}.
\end{itemize}
\end{document}
